% ECONOMICS_PROBLEM: {"id": "C90-284", "title": "Adaptive experimental inference", "jel_code": "C90", "source_id": "OP-0428"}
% !TeX program = xelatex
\documentclass[11pt,a4paper]{article}
\usepackage[margin=23mm]{geometry}
\usepackage{fontspec}
\setmainfont{Latin Modern Roman}
\usepackage{amsmath,amssymb,mathtools}
\usepackage{xcolor,enumitem,booktabs,longtable,array,xurl,hyperref}
\definecolor{muted}{HTML}{53636C}
\hypersetup{colorlinks=true,urlcolor=blue,linkcolor=blue}
\setlength{\parindent}{0pt}
\setlength{\parskip}{5pt}
\newcommand{\R}{\mathbb R}
\newcommand{\E}{\mathbb E}
\newcommand{\N}{\mathbb N}
\newcommand{\ind}{\mathbf 1}
\newcommand{\Var}{\operatorname{Var}}
\newcommand{\Cov}{\operatorname{Cov}}
\newcommand{\supp}{\operatorname{supp}}
\DeclareMathOperator*{\argmin}{arg\,min}
\DeclareMathOperator*{\argmax}{arg\,max}
\newcommand{\entrymeta}[3]{{\sffamily\small\color{muted}\textbf{\texttt{#1}}\quad|\quad #2\par #3\par}}
\newcommand{\Problemref}[1]{\texttt{#1}}
\newcommand{\JELref}[2]{\texttt{#2} (JEL \texttt{#1})}
\begin{document}
\section*{Adaptive experimental inference}
\textbf{Problem ID:} \texttt{C90-284}\quad \textbf{JEL:} \texttt{C90}\par
% BEGIN ECONOMICS STATEMENT
\label{problem:OP-0428}
{\sffamily\small\color{muted}Original subject group: Econometrics, causal inference and measurement\par}
\label{definitions:problem:30007710}
\textbf{Audit classification:} Background; proposed extension. Correct authors verified; published2025 DOI added. Adaptive inference already has solutions, and a bounded IPS martingale gives simultaneous coverage here even under unrestricted conditional drift. Only the explicitly specified optimization frontier is a proposed research task.
\subsection*{Definitions and precise research target}
Consider a sequential experiment over \(T\) participants with two treatments \(a\in\{0,1\}\) and bounded potential outcomes \(Y_t(a)\in[0,1]\). Conditional mean rewards \(\mu_{a,t}=E[Y_t(a)\mid\mathcal H_{t-1}]\) may change with time, where \(\mathcal H_{t-1}\) is recorded history. Assume conditional randomization given history and an almost-sure drift budget \(\sum_{t=2}^{T}\max_a|\mu_{a,t}-\mu_{a,t-1}|\leq V\). Assignment probability \(p_t\), measurable from history, satisfies \(\epsilon\leq p_t\leq1-\epsilon\) for fixed \(0<\epsilon<1/2\). Define the running average effect \(\tau_t=t^{-1}\sum_{s=1}^{t}(\mu_{1,s}-\mu_{0,s})\) and welfare regret \(R_T=\sum_{t=1}^{T}\{\max_a\mu_{a,t}-\mu_{A_t,t}\}\), where \(A_t\) is the assigned treatment. Develop assignment and inference rules whose confidence intervals \(C_t\) satisfy \(\Pr(\tau_t\in C_t\text{ for every }t\leq T)\geq1-\alpha\), for stated \(\alpha\in(0,1)\), uniformly over the bounded-outcome drift class. Characterize attainable trade-offs between expected regret and interval width as functions of \(T,V,\epsilon\), and compare them with fixed randomization. Regret compares one-step choices along realized history, rather than alternative state trajectories. Fixed positive exploration can imply linear regret even without drift. Arbitrary unbounded drift is excluded explicitly. Adaptive inference already has established solutions; the nonstationary joint welfare-and-inference optimization is a proposed extension, and the supporting reference does not verify its exact open status. Outcomes are assessed in a prespecified participant population with consistent measurement across the observation horizon.

\textbf{Definitions, assumptions and mathematical qualifications.} Let \(\mathcal H_{t-1}\) contain all prior assignments/outcomes. Require \(A_t\perp(Y_t(0),Y_t(1))\mid\mathcal H_{t-1}\), with \(p_t=\Pr(A_t=1\mid\mathcal H_{t-1})\). Observed \(Y_t=Y_t(A_t)\) yields score \(X_t=A_tY_t/p_t-(1-A_t)Y_t/(1-p_t)\), whose conditional mean is \(\mu_{1,t}-\mu_{0,t}\). For fixed horizon \(T\), Hoeffding's martingale bound and a union bound give valid intervals centered at \(t^{-1}\sum_{s\leq t}X_s\), intersected with \([-1,1]\), with radius \(\epsilon^{-1}\sqrt{2\log(2T/\alpha)/t}\). Thus simultaneous coverage exists without the drift restriction; existence is not open. Study the attainable frontier of \((\sup_M E_MR_T,\sup_M E_M\sum_t|C_t|)\) over a specified outcome/drift class and admissible designs. Fixed exploration yields \(E R_T\geq\epsilon\Delta T\) when stationary arm means differ by \(\Delta>0\).
\subsection*{Variants and assumption changes}
\begin{enumerate}
\item \textbf{Editorial, fixed exploration benchmark.} Keep \(\epsilon>0\); compare the displayed valid intervals and unavoidable linear regret with stronger variance-adaptive inference or tighter Pareto bounds.
\item \textbf{Editorial, declining exploration.} Permit \(\epsilon_t\downarrow0\) according to a stated rule and impose a drift budget \(V\); characterize how increasing inverse-probability variance trades against dynamic regret and anytime interval length.
\end{enumerate}
\subsection*{References and source audit}
\textbf{Checked source.} arXiv:2405.01281v1 (2 May 2024), Sections 2--5; Annual Review of Statistics and Its Application 12 (2025), pp. 407--423. The survey describes existing valid-inference methods; it does not declare this nonstationary frontier open. Primary preprint text and published metadata checked. Preprint cover misspells Demistifying; published title is Demystifying. Published DOI 10.1146/annurev-statistics-040522-015431. \url{https://arxiv.org/pdf/2405.01281}.
\begin{enumerate}\raggedright
\item\hypertarget{definitions:source:30007710:1}{} Aurélien Bibaut; Nathan Kallus. 2024. \emph{Demystifying Inference After Adaptive Experiments}. arXiv:2405.01281v1 (2 May 2024), Sections 2--5; Annual Review of Statistics and Its Application 12 (2025), pp. 407--423.
\url{https://arxiv.org/pdf/2405.01281}.
DOI: \url{https://doi.org/10.48550/arXiv.2405.01281}.
\end{enumerate}

\subsection*{Common conventions: Source mathematical conventions}

These conventions apply to each entry unless explicitly replaced.

\textbf{Models and identification.} A primitive model \(m\) specifies all probability laws, feasible choices, information, equilibrium and measurement rules used by its target. The admissible class \(\mathcal M\) is fixed independently of outcomes. If \(P_O(m)\) is its observable probability law and \(\tau(m)\) the target, its identified set at observable law \(P\) is
\[
 \mathcal I_\tau(P)=\{\tau(m):m\in\mathcal M,\ P_O(m)=P\}.
\]
Point identification means that this set is a singleton. An empty set signals incompatibility of data and assumptions. Coordinate bounds need not describe the joint set. Bounds are sharp only if every retained target is attainable or arbitrarily approachable by compatible models. Finite bounds require explicit restrictions on unobserved quantities; unrestricted confounding, hidden wealth, extrapolation or arbitrary equilibrium selection can yield uninformative sets.

\textbf{Causal interventions.} A potential outcome \(Y_i(p)\) is the outcome under a complete policy regime \(p\), including timing, financing, information and market spillovers. The target population and units are fixed before treatment. With interference, use the complete assignment vector or a justified exposure mapping; individual no-interference notation alone cannot identify an economy-wide intervention. A difference of mean potential outcomes does not require the cross-world joint distribution. Conditioning on post-treatment migration, survival, enrollment or employment changes the estimand and needs separate justification.

\textbf{Statistical statements.} An estimator, confidence set, forecast or test specifies a sampling law, model class and loss/coverage criterion. Uniform \((1-\alpha)\)-coverage means \(\inf_{m\in\mathcal M}\Pr_m\{\tau(m)\in C_n\}\geq1-\alpha\), or an explicitly asymptotic version. The whole parameter space trivially has coverage, so informativeness requires a stated width, risk or power objective. A forecasting improvement is relative to a named benchmark, information set and evaluation population.

\textbf{Equilibrium and optimization.} An equilibrium comprises feasible plans, optimality, beliefs where needed, budget constraints and all market/resource clearing. The primitives or a stated benchmark define each correspondence \(\mathcal E(p;m)\). Nonemptiness is assumed or proved, never inferred just from compact policy sets. A selected-equilibrium target uses a declared selection rule; a robust target quantifies over all permitted equilibria. Use suprema or infima unless attainment is established. An \(\varepsilon\)-optimal policy is feasible and within \(\varepsilon\) of the finite optimum value. Report an empty feasible set separately.

\textbf{Welfare and accounting.} Normative utility, interpersonal normalization, social weights, numeraire, discounting and the use of public receipts are primitives. Behavioral utility need not coincide with the welfare criterion. Household welfare includes ownership income and rents once; adding profits again to compensating variation generally double-counts them. A monetary equivalent is defined by an explicit transfer and price convention, with monotonicity and range conditions. Average effects, mechanism contributions, transfers and welfare losses are different targets. Infinite sums require convergence and dynamic feasible plans require terminal or no-Ponzi conditions.

\textbf{Source versions.} A working-paper date, revision date and journal publication year are distinguished. A DOI for a journal version is not automatically the DOI of an earlier manuscript. Locators refer to the stated version and printed pages where specified. A metadata-only check does not verify a claim about a full-text conclusion. Every entry's source subsection narrows the provenance claim accordingly.
% END ECONOMICS STATEMENT
\end{document}
